% Generated draft: requires mathematical and source review.
\documentclass[11pt]{article}
\usepackage[margin=25mm]{geometry}
\usepackage{fontspec}
\setmainfont{Latin Modern Roman}
\usepackage{amsmath,amssymb,mathtools,unicode-math}
\setmathfont{Latin Modern Math}
\usepackage{enumitem,needspace,xurl,hyperref}
\setlength{\parindent}{0pt}
\setlength{\parskip}{5pt}
\setlength{\emergencystretch}{3em}
\newcommand{\N}{\mathbb{N}}
\newcommand{\Z}{\mathbb{Z}}
\newcommand{\Q}{\mathbb{Q}}
\newcommand{\R}{\mathbb{R}}
\newcommand{\C}{\mathbb{C}}
\newcommand{\F}{\mathbb{F}}
\newcommand{\A}{\mathbb{A}}
\newcommand{\PP}{\mathbb{P}}
\newcommand{\E}{\mathbb{E}}
\newcommand{\eps}{\varepsilon}
\newcommand{\dd}{\,\mathrm{d}}
\begin{document}
\section*{UnsolvedMath ID: 30006258}
\textbf{Source problem code:} catalogue entry
\subsection{Definitions and mathematical statement}
Write \(A=u_q(\mathfrak g)\) for the small quantum group intended by the supplied record, with \(\mathfrak g\) of type \(A\), and write
\[
Z=z(A)=\{z\in A:za=az\text{ for every }a\in A\}.
\]
For a specified action of the intended group \(G\) on \(Z\), its invariant subalgebra is
\[
Z^G=\{z\in Z:g\cdot z=z\text{ for every }g\in G\}.
\]
The record does not specify the ground field, the rank, the relation between \(G\) and \(\mathfrak g\), the order \(\ell\) of the root of unity \(q\), the admissibility restrictions on \(\ell\), or the action defining these invariants. These data must be fixed from the source before this becomes a fully quantified problem.

\textbf{Question recorded in the supplied statement.} For the intended admissible small quantum groups of type \(A\), is there an isomorphism of algebras
\[
z\bigl(u_q(\mathfrak g)\bigr)^G\cong z\bigl(u_q(\mathfrak g)\bigr)?
\]
The record does not specify whether additional compatibility is required of this isomorphism or whether the intended comparison instead concerns particular blocks of the center.

The supplied statement also asks for the bigraded and multiplicative structure of \(Z\), and for a natural action of the intended Weyl group \(W\) compatible with a diagonal-coinvariant model. The bigrading, the precise diagonal-coinvariant algebra, the comparison map or isomorphism, the blocks to which it applies, and the meaning of compatibility are not supplied. Consequently, this additional request cannot yet be expressed as a uniquely determined graded-algebra or representation-theoretic assertion.

\textbf{Qualification from the original excerpt.} The original excerpt does not explicitly state the displayed question. It reports that an algebra isomorphism between the center and certain cohomology of an affine Springer fiber was conjectured in reference [3] and subsequently proved. It then attributes a partial confirmation of Conjecture 3 to references [5] and [3], subject to unspecified restrictions on \(\ell\). Its Theorem 4 involves an integral weight \(\lambda\in P\) and a parabolic subgroup, also denoted \(P\), whose Weyl group \(W_P\subset W\) stabilizes \(\lambda\). The action used for this stabilizer is unspecified. The ensuing dimension formula is corrupted in the supplied extraction, so neither its invariant superscripts nor its block and diagonal-coinvariant notation can be recovered reliably. No repaired theorem or identification with the displayed question is asserted here.
\subsection{Short English statement}
The supplied formulation asks whether the invariant center of a type \(A\) small quantum group is isomorphic to its full center, and seeks the center's bigrading, multiplication, and a compatible Weyl-group action. Essential parameters and comparison structures are missing, and the original excerpt supplies a corrupted partial-result statement rather than the explicit question. This is a statement-only draft requiring human mathematical review.
\subsection{Sources}
\begin{itemize}
\item Supplied local registry record: permanent UnsolvedMath ID \(30006258\), title \emph{Invariant Centers of Small Quantum Groups}, exact source problem code \texttt{catalogue entry}, in the Oberwolfach Reports Open Problems collection. The record assigns the year 2025 and names no individual proposer.
\item Supplied report citation: \emph{Algebraic Groups} (2025), Oberwolfach Reports, pp.\ 941--1010, DOI \texttt{10.4171/owr/2025/20}. Supplied source URL: \url{https://doi.org/10.4171/owr/2025/20}. No narrower page location for this problem is supplied.
\item Supplied literature source, labeled ``OWR 2025 report'': \url{https://ems.press/content/serial-article-files/51849}.
\item Supplied literature source, labeled ``Lachowska--You, singular blocks'': \url{https://doi.org/10.1112/plms.12193}. The record does not identify this citation with either numbered reference [3] or [5].
\item Supplied literature assessment dated \texttt{2026-08-22}: the record labels the problem open, says that important questions remain unresolved, and reports that a nearby invariant-center conjecture is known only in small type-\(A\) rank and fails in several other types. It therefore records no general resolution. This is supplied dated curation derived from the OWR material, not an independently verified assessment; its time-sensitive status and cited claims require independent verification. The ``nearby'' conjecture is not identified precisely enough to equate it with the displayed question.
\item The registry marks the statement \texttt{corrected\_verified} and reports that the official report was checked and its notation repaired. These are supplied provenance claims, not verification performed for this draft. The original excerpt's references [3], [5], and Conjecture 3 are not supplied in full.
\end{itemize}
\end{document}
